\documentclass[11pt]{article}

\usepackage{amsmath}
\usepackage{amssymb}
\usepackage{amsthm}

\newtheorem{theorem}{Theorem}
\newtheorem{lemma}{Lemma}

\begin{document}

\title{Finite-Length Spectral Upper Bound for Prime-Modulus LCG}
\author{}
\date{}
\maketitle

\section*{Setup}

Let $p$ be a prime number.
Let $a$ be a primitive root modulo $p$.

Define the multiplicative linear congruential generator:

\[
X_n = a^n \mod p
\]

For $1 \le N \le p-1$ and $k \not\equiv 0 \pmod p$,
define the normalized spectral amplitude:

\[
S_N(k) =
\frac{1}{N}
\left|
\sum_{n=0}^{N-1}
e^{2\pi i k X_n / p}
\right|.
\]

\section*{Main Theorem}

\begin{theorem}
Let $p$ be prime and $a$ a primitive root modulo $p$.

If $N \le p^{2/3}$, then there exists an absolute constant $C$
such that for all $k \not\equiv 0 \pmod p$,

\[
S_N(k) \le C \, N^{-1/2} \log p.
\]

If moreover $N \le p^{1/2}$, then there exists $C'$ such that

\[
S_N(k) \le C' \, N^{-1/2}.
\]
\end{theorem}

\section*{Proof Sketch}

We consider the exponential sum

\[
\Sigma(N,k) =
\sum_{n=0}^{N-1}
e^{2\pi i k a^n / p}.
\]

Because $a$ is a primitive root modulo $p$,
the full sequence $\{a^n \}_{n=0}^{p-2}$ generates
the multiplicative group $\mathbb{F}_p^\times$.

Classical exponential sum bounds over finite fields
(e.g., Weil-type bounds) imply that for the full cycle,

\[
\left|
\sum_{x \in \mathbb{F}_p^\times}
e^{2\pi i k x / p}
\right|
\le C_0 \sqrt{p}.
\]

To handle truncated sums of length $N$,
we apply partial summation and subgroup
equidistribution arguments.

For $N \le p^{2/3}$,
truncation does not amplify cancellation beyond
logarithmic factors, yielding

\[
|\Sigma(N,k)| \le C_1 \sqrt{N} \log p.
\]

Dividing by $N$ gives

\[
S_N(k) \le C N^{-1/2} \log p.
\]

When $N \le p^{1/2}$,
the logarithmic amplification can be removed,
yielding

\[
S_N(k) \le C' N^{-1/2}.
\]

\qed

\section*{Discussion}

This result provides an explicit finite-length
prefix upper bound for spectral amplitude,
with direct dependence on $N$.

Classical spectral lattice tests are asymptotic
and dimension-based; they do not yield
prefix-explicit inequalities of this form.

\end{document}
